% ============================================================
% PAPER B — The Obstruction Spectrum of Weyl Context Families
%           and the 2-adic Tower
% arXiv submission source.

\documentclass[11pt]{article}
\usepackage[margin=1in]{geometry}
\usepackage{amsmath,amssymb,amsthm,mathtools}
\usepackage[hidelinks]{hyperref}
\usepackage{lmodern}

\newtheorem{theorem}{Theorem}
\newtheorem{lemma}[theorem]{Lemma}
\newtheorem{proposition}[theorem]{Proposition}
\newtheorem{corollary}[theorem]{Corollary}
\newtheorem{conjecture}[theorem]{Conjecture}
\theoremstyle{remark}
\newtheorem{remark}[theorem]{Remark}

\newcommand{\Z}{\mathbb{Z}}
\newcommand{\F}{\mathbb{F}}
\newcommand{\ip}[2]{\langle #1,#2\rangle}
\newcommand{\ipz}[2]{\langle #1,#2\rangle_{\Z}}
\DeclareMathOperator{\col}{col}
\DeclareMathOperator{\rowspace}{rowspace}
\DeclareMathOperator{\supp}{supp}

\title{The Obstruction Spectrum of Weyl Context Families\\ and the 2-adic Tower}
\author{Manuel Flores Gordillo\thanks{ORCID \texttt{0009-0006-8246-0343};
\texttt{manuel.flores@columbia.edu}. Verification scripts and data:
\url{https://github.com/manuflog/contextuality-obstructions}.}}
\date{July 2026}

\begin{document}
\maketitle

\begin{abstract}
State-independent contextuality of qudit stabilizer measurements is not just present or absent --- it
comes in a definite strength. For $m$-qudit Weyl--Heisenberg context families in local dimension $d$,
the achievable all-versus-nothing obstruction value is exactly $\{0,d/2\}$ for every even $d$ and
$\{0\}$ for odd $d$. \textbf{This value classification --- the even/odd dichotomy and the fixed value
$d/2$ --- is due to Abramsky, Cercelescu and Constantin~\cite{ACC24}}, whose parity obstruction
$2k\equiv 0\pmod d$ (their Thm.~16) and normal-form classification (their Thm.~21) settle it; our
Obstruction Spectrum Theorem (Theorem~\ref{thm:spectrum}) recovers their $2S\equiv 0\pmod d$ in
explicit Weyl-certificate language. On this basis we develop three further layers. (i) A closed-form
calculus: contextuality is decided by an integer commutator carry
$b(C)=\sum_{\{u,v\}}\ipz{u}{v}/d \pmod 2$ (Theorem~\ref{thm:carry}) --- no phases or matrices --- with
an exact value-bit formula (Theorem~\ref{thm:valuebit}) and a cycle criterion
(Theorem~\ref{thm:cycle}), in which the carry invariant $Q$ has the algebraic form of a
Pontryagin-square quadratic refinement (we do not construct the cohomology operation here).
(ii) A $2$-adic tower structure: the obstruction lives at the top layer and no single lift carries it
(Depth Rigidity, Theorem~\ref{thm:depth}), with a fiber doubling law governing its propagation
$d\to 2d$ (generativity obstruction $G=\lfloor\Sigma M/2\rfloor\oplus\lfloor K/2\rfloor$).
(iii) A three-device superconducting replication of the $d=2$ witness. The obstruction's geometric
home --- a canonical degree-$2$ class of order $n$ on the stack $[X/\mathrm{PU}(n)]$ of maximal
abelian subalgebras --- is developed in the companion paper.
\end{abstract}

\section{Setting and the question}
Fix $m$ qudits of local dimension $d$ and symmetric Weyl lifts
$W(v)=\tau^{-q(v)}X^{a_1}Z^{b_1}\otimes\cdots$, $\tau=e^{i\pi/d}$, integer representatives $0..d-1$,
$q(v)=\sum_i a_ib_i$. A \emph{context} $C$ is a word of pairwise-commuting $v$'s
($\ip{u}{v}\equiv 0 \bmod d$) with $\sum v\equiv 0 \bmod d$; its product is a scalar $\omega^{s(C)}$,
$\omega=e^{2\pi i/d}$. A \emph{certificate} is $\lambda\in\Z_d^{\mathrm{ctx}}$ with
$\lambda^{\mathsf T}A\equiv 0\pmod d$, where $A$ is the context-by-observable multiplicity matrix; its
value is $S=\lambda^{\mathsf T}s\bmod d$. A nonzero value is a state-independent all-versus-nothing
(AvN) obstruction: no $\Z_d$-valued noncontextual assignment is consistent with the family. Which values $S$ can occur was settled by~\cite{ACC24}; our aim is to recast that answer in explicit
certificate language and add the carry calculus, the $2$-adic tower, and a hardware replication. Let
$H(d,m)\subseteq\Z_d$ be the set of achievable values.

\paragraph{Practical takeaway.} Three facts are directly usable. (i) \emph{Even/odd dichotomy}:
state-independent Weyl contextuality exists exactly for even $d$; in odd $d$ every such family is
noncontextual ($H=\{0\}$), so an experiment seeking a preparation-untrusted contextuality witness must
use even local dimension. (ii) \emph{Fixed strength}: when present, the obstruction takes the single
value $d/2$ (the sign $-1=\omega^{d/2}$) --- there is no continuum of strengths to tune, and
correspondingly no additive resource monotone (the obstruction is a saturating support, not a charge).
(iii) \emph{Cheap witnesses}: contextuality is decided by an integer commutator carry
$b(C)=\sum_{i<j}\ip{c_i}{c_j}/d \bmod 2$ with no matrices or phases, and a witness at $2d$ can be
generated from one at $d$ by the fiber construction (though genuine AvN minimality is a $\Z_d$-cycle
condition, Proposition~\ref{prop:zd}, not the $\Z_2$ shadow). The natural application is
state-independent certification of quantumness in even-$d$ hardware: the witnesses hold for any input
state, including the maximally mixed one, so they benchmark a device without trusting state
preparation. We claim no bridge to code performance or $T$-count; the obstruction value classifies a
cohomological invariant, not a computational resource.

\paragraph{Operational certification (application).} The dichotomy gives a concrete measurement
protocol. For a Weyl family one measures, in each context $C$, the joint outcomes of its commuting
generators and forms the product statistic; the noncontextual functional
$S=\sum_C \cos\!\big(2\pi(\hat s_C-s(C))/d\big)$ has a classical bound below the quantum value exactly
when $H\neq\{0\}$, i.e.\ for even $d$. Two features make this device-friendly. First, the context
product operator is the scalar $\omega^{s(C)}I$, so the statistic is fixed by the operator algebra and
is independent of the input state (the maximally mixed state included); state-preparation error does
not degrade it, and only measurement imperfection enters. A protocol caveat is essential here and is
easily overlooked: the scalar identity $\prod_{O\in C}O=\omega^{s(C)}I$ means that a single joint
measurement of a context---diagonalizing the commuting family and reading all outcomes from one basis
label---returns the product $\omega^{s(C)}$ on every shot as an algebraic identity, independently of
the device, so it certifies nothing (it reproduces the ideal value even on random outcomes). The
measurement-limited character is recovered only when the observables of a context are measured in
physically separate operations (sequential non-destructive readout, or one ancilla per observable):
the outcomes are then no longer bound by the identity, and measurement error genuinely degrades the
correlator below $|\omega^{s(C)}|=1$. This preparation-untrusted, measurement-limited character of
state-independent contextuality is well established experimentally on trapped ions, photons, and
superconducting qubits~\cite{Kirchmair,Qu,EJP}; we use it rather than claim it. Second, and specific
to the present classification, the quantity being certified is not merely the presence of
contextuality but its \emph{value}: the graded invariant $S=d/2$, together with the even/odd signature
(the same construction is provably noncontextual in odd $d$), is a structural fingerprint of the
obstruction spectrum. A ququart-encoded ($d=4$) witness realizing this on qubit hardware needs only
shallow circuits and tolerates measurement error at the several-percent scale, well within current
devices; the value-graded hardware certification is developed separately.

\section{The Obstruction Spectrum Theorem}
\begin{theorem}[Obstruction Spectrum]\label{thm:spectrum}
For every even $d$, every $m$, every Weyl context family and every certificate, $2S\equiv 0\pmod d$;
hence $S\in\{0,d/2\}$. For odd $d$ (in the Gross convention $W'(v)=\omega^{-2^{-1}q(v)}T(v)$), $S=0$
always. Thus the only state-independent obstruction value achievable by Weyl context families in any
even dimension is $-1=\omega^{d/2}$.
\end{theorem}
\begin{proof}
The Weyl product obeys $W(u)W(v)=\tau^{c(u,v)}W((u+v)\bmod d)$ with
$c(u,v)=-q(u)-q(v)+2\beta(u,v)+q((u+v)\bmod d)$, $\beta(u,v)=\sum_i u_{b,i}v_{a,i}$ (machine-verified
against matrices at $d=4,6,9,16$, $25/25$ contexts each). Telescoping a context word from the
identity, the $q$-terms of the partial sums cancel:
\begin{equation}\label{eq:mod2d}
2s(C)\equiv -\sum_i q(u_i)+2\sum_{j<i}\beta(u_j,u_i)\pmod{2d},
\end{equation}
$B(C)=\sum_i \beta(P_i,u_i)$, $P_i$ the prefix sums; reducing mod $d$, bilinearity gives
$B(C)\equiv\sum_{j<i}\beta(u_j,u_i)$. On commuting pairs $\beta(u,v)-\beta(v,u)=\ip{u}{v}\equiv 0$, so
$\beta$ is symmetric mod $d$ there; polarizing $\beta(\sigma,\sigma)$ at $\sigma=\sum_i u_i\equiv 0$,
using $\beta(u,u)=q(u)$, yields $2\sum_{j<i}\beta(u_j,u_i)\equiv -\sum_i q(u_i)\pmod d$. Substituting
into~\eqref{eq:mod2d} reduced mod $d$:
\[
2s(C)\equiv -2\sum_{u\in C}q(u)\pmod d :
\]
the doubled phase is a coboundary with explicit potential $-q$. Therefore
$2S=\lambda\cdot 2s\equiv -2(\lambda^{\mathsf T}A)\cdot q\equiv 0\pmod d$. For odd $d$, $2$ is
invertible mod $d$ and the Gross-convention telescoping gives $s(C)\equiv-\sum_{u\in C}q(u)\pmod d$
outright, so every certificate value vanishes.
\end{proof}

\noindent\emph{Verification.} The potential identity holds $0/3000$ at each of $d=2,4,6,8,12,16$ and
$0/18521$ across mixed $(d,m)$; the mod-$2d$ strengthening fails ${\sim}50\%$ --- the residue is
exactly the obstruction bit, so the theorem is sharp. Odd-$d$ triviality: single value $\{0\}$ at
$d=3,5,9$ (\texttt{spectrum\_test2.py}; full ledger in the repository \texttt{INDEX.md}).

\begin{corollary}[The obstruction is one bit]\label{cor:bit}
$s\bmod (d/2)$ is exact with potential $-q$; all obstruction content lives in
$b(C):=\big(s(C)+\sum_{u\in C}q(u)\big)/(d/2)\in\Z_2$, and $S=(d/2)\,(\bar\lambda\cdot b\bmod 2)$ with
$\bar\lambda=\lambda\bmod 2$. Only $\lambda\bmod 2$ carries value; any attaining certificate has odd
entries. The value group is always $\subseteq\mu_2\subset\mu_d$: the physics factors through the real
shadow at every even $d$, tying the two-tier structure of Paper~A to a quantitative statement here.
\end{corollary}

\begin{remark}[No higher-order qudit parity proofs]\label{rem:nohigher}
Theorem~\ref{thm:spectrum} fixes the value in $\{0,d/2\}$ for any incidence structure, of any
multiplicity: if each observable lies in $r$ contexts, the concatenated word has every letter $r$
times, yet the achievable value is still only $0$ or $d/2$, independent of $r$. Since $d/2$ is the
order-two element of $\mu_d$ (Corollary~\ref{cor:bit}), there is no obstruction of order $>2$ in any
Weyl family --- the value is always $2$-primary and visible to the mod-$2$ shadow. In particular the
natural expectation that high-multiplicity covers ($d\mid r$, $r>2$) would yield genuinely order-$d$
``qudit parity proofs'' fails: any such cover, if contextual, still carries the single bit $d/2$.
(Verified: stacking a $d=4$ certificate to multiplicities $r=4,6$ keeps the value group
$\subseteq\{0,2\}$, never $\{1,3\}$.)
\end{remark}

\section{The commutator-carry form}
\begin{theorem}[Carry form]\label{thm:carry}
Exactly, $2s(C)\equiv -\sum_i q(u_i)+2\sum_{j<i}\beta(u_j,u_i)\pmod{2d}$, and consequently
\[
b(C)\equiv\sum_{\{u,v\}\subseteq C}\ipz{u}{v}/d\pmod 2
\]
--- the obstruction bit is the total commutator carry. The whole theory is combinatorial: no phases,
no matrices, only symplectic carries.
\end{theorem}
\noindent\emph{Verification.} Carry form $0/2500$ at each $d=2,4,6,8,12,16$, and $0/800$ per $d$ with
zero letters and high multiplicity.

\section{Depth Rigidity}
A lift of a level-$d$ family to level $2d$ chooses $x_v\in\{0,1\}^{2m}$ per observable
($\tilde v=v+d\,x_v$) keeping all contexts commuting and sum-zero mod $2d$ (a $\mathrm{GF}(2)$-linear
system with carry right-hand sides).
\begin{theorem}[Depth Rigidity]\label{thm:depth}
(a) If the commutation block of the lift system is solvable then
$b(C)\equiv\sum_{v\in C}\ip{x_v}{v}\pmod 2$ is a coboundary, so every certificate value of the lifted
family is $0$. (b) Conversely, if a level-$d$ family attains (some $\lambda$ with value $d/2$), the
$\bar\lambda$-weighted commutation equations have identically zero left side and right side
$\sum_j\bar\lambda_j b(C_j)\equiv 1$; hence attaining families never lift. The obstruction sits at the
top $2$-adic layer and is carried by no single section.
\end{theorem}
\noindent\emph{Verification.} On the $d=8$ pool the full mod-$8$ kernel has value group $\{0,4\}$ but
the $2$-divisible / liftable part gives only $\{0\}$; the $d=8$ certificate $\lambda$ is odd on
$62/89$ contexts. The dual no-lift witness for the $d=8\to 16$ step is machine-verified (data in the
repository).

\section{The value-bit formula}
\begin{theorem}[Value-bit formula]\label{thm:valuebit}
For any mod-$2d$ kernel element $\tilde\lambda$ of a pure fiber pool, with odd support $\supp$ and
lift data $u=\tilde v\bmod d$, $x=\tilde v\ \mathrm{div}\ d$, the value bit $P$ ($S=d\,P$) is
\[
P\equiv \tfrac12\Big[\textstyle\sum_{\supp}K+\sum_{\supp}\Lambda(x)+(d/2)\sum_{\supp}Q(x)\Big]\pmod 2,
\]
$K=\sum_{i<i'}\ipz{u_i}{u_{i'}}/d$, $\Lambda(x)=\ip{\sum_i x_i}{\sigma}-\sum_i\ip{x_i}{u_i}
-2\sum_i\ip{x_i}{P_i}$, $Q(x)=\sum_{i<i'}\ip{x_i}{x_{i'}}$; the bracket is always even. The
$(d/2)$-terms are visible mod $2$ only for $d\equiv 2\pmod 4$, reproducing the regime split.
\end{theorem}
\noindent\emph{Verification.} $0/550$ on each of $\mathrm{PM}\to 4$, $\mathrm{cert4}\to 8$, and
non-attaining pools, non-vacuously (\texttt{Wformula.py}).

\section{Attainment and the classification}
The upper bound (Theorem~\ref{thm:spectrum}) is complemented by attainment at every even $d$.
\begin{theorem}[Attainment / classification]\label{thm:attain}
For every even $d$ and $m\ge 2$, the value $d/2$ is attained, so $H(d,m)=\{0,d/2\}$; for odd $d$,
$H(d,m)=\{0\}$.
\end{theorem}
\noindent For general even $d$ the lower bound is the ROZF construction~\cite{ROZF}, whose even-$d$
attainment we translate to Weyl certificates; the \emph{generality} of attainment is inherited from
that construction, while the explicit, self-contained certificates below --- independently
matrix-verified at $d=8$ (value $4$, $89$ contexts) and $d=16$ (value $8$, $363$ contexts) --- are
new. The ROZF face construction attains $\{0,d/2\}$ at $d=4,6,8,10,12,16,20$; \texttt{verify\_cert8.py}
confirms the $d=8$, $S=4$ certificate on raw $64\times 64$ matrices (commuting, scalar,
$\lambda^{\mathsf T}A\equiv 0\bmod 8$, value $4$, system provably infeasible), and
\texttt{verify\_cert16.py} confirms $d=16$, $S=8$ on raw $256\times 256$ matrices. The classification
is thus a theorem, not contingent on the tower results of the next section; the $d=16$ certificate in
particular shows there is no depth ceiling.

\section{A closed-form contextuality criterion}
The value-bit machinery yields, for even $d$, an exact and elementary criterion for contextuality in
terms of a single symplectic invariant of each cycle. Fix a family of contexts $\{C_r\}$, each a set
of Weyl observables (symplectic vectors in $\Z_d^{2m}$) summing to $0\bmod d$; let $A$ be its $\Z_2$
context--observable incidence and $b_r=\sum_{i<j\in C_r}\ip{c_i}{c_j}/d\bmod 2$ the carry vector. A
\emph{cycle} is a left-kernel vector $\lambda$ ($\lambda^{\mathsf T}A\equiv 0$), i.e.\ an
even-multiplicity combination of contexts.
\begin{theorem}[Cycle value $=$ symplectic self-pairing]\label{thm:cycle}
Let $d$ be even and $\lambda$ a cycle. Then
\[
\lambda\cdot b=Q(\lambda):=\sum_{a<b\ \text{in the observable multiset of }\lambda}\ip{v_a}{v_b}/d
\pmod 2.
\]
Consequently the family is contextual if and only if some cycle $\lambda$ has $Q(\lambda)$ odd.
\end{theorem}
\begin{proof}
By definition $\lambda\cdot b=\sum_{r:\lambda_r=1}\sum_{i<j\in C_r}\ip{c_i}{c_j}/d$, the intra-context
pairs. Splitting $Q$ into intra- and inter-context pairs, $Q=\text{intra}+\text{inter}$ with
$\text{inter}=\sum_{a<b}\sum_{u\in C_a,v\in C_b}\ip{u}{v}/d$. By bilinearity
$\sum_{u\in C_a,v\in C_b}\ip{u}{v}=\ip{S_a}{S_b}$ where $S_r\in\Z^{2m}$ is the integer sum of the
observables of $C_r$. Each context sums to $0\bmod d$, so $S_r=d\,t_r$ for an integer vector $t_r$;
hence $\ip{S_a}{S_b}=d^2\ip{t_a}{t_b}$ and $\text{inter}=\sum_{a<b}d\ip{t_a}{t_b}\equiv 0\pmod 2$ since
$d$ is even. Thus $Q=\text{intra}=\lambda\cdot b$. The family is contextual iff
$b\notin\rowspace(A)$, i.e.\ iff some left-kernel $\lambda$ has $\lambda\cdot b=1$, i.e.\ $Q(\lambda)$
odd.
\end{proof}
The invariant $Q$ has the algebraic form of a Pontryagin-square quadratic refinement of the cycle:
contextuality is the oddness of this quantity on some cycle. This is the sense in which the value
obstruction resembles the support of a type-III (Pontryagin) anomaly of the Weyl--Heisenberg group ---
the obstruction saturates rather than adds under stacking. We do not here construct the cochain
complex, coefficient sequence, and cohomology operation that would identify $Q$ with the Pontryagin
square on the nose; that precise identification is left open (see \S\ref{sec:open}). Generic cycles have $Q=0$ (verified $300/300$ on random families); attaining
certificates are precisely those carrying a cycle with $Q=1$. For odd $d$ the halving step fails and
the family is noncontextual ($H=\{0\}$), consistent with the absence of odd-$Q$ cycles.

\begin{proposition}[The exact $\Z_d$ criterion; $Q$ is its shadow]\label{prop:zd}
Let $M$ be the $\Z_d$ incidence matrix and $\gamma(C)=\sum_{i<j}\ip{c_i}{c_j}/d\in\Z_d$ the context
scalar. By the Fredholm alternative over $\Z_d$, a family is state-independent contextual
($\gamma\notin\col(M)$) if and only if there is a $\Z_d$-cycle $\lambda$
($\lambda^{\mathsf T}M\equiv 0\pmod d$) with $\lambda\cdot\gamma\not\equiv 0\pmod d$. The odd-$Q$
criterion of Theorem~\ref{thm:cycle} is the mod-$2$ shadow of this. For \emph{even} $d$ the shadow
is sound, and the proof is two lines: if $Mx\equiv\gamma\pmod d$ then also $Mx\equiv\gamma\pmod 2$,
so any $\lambda$ with $\lambda^{\mathsf T}M\equiv 0\pmod 2$ has
$\lambda\cdot\gamma\equiv\lambda^{\mathsf T}Mx\equiv 0\pmod 2$. Hence the shadow can \emph{never}
fire on a noncontextual system: odd-$Q$ has no false positives at any even $d$.

It is, however, \emph{incomplete}, and in the opposite direction: a $\Z_d$-cycle certifies
contextuality when $\lambda\cdot\gamma\not\equiv0\pmod d$, and if that value is nonzero but
\emph{even} its mod-$2$ reduction vanishes, so the shadow stays silent on a genuinely contextual
system. An explicit non-degenerate $d=8$ witness is pinned in
\texttt{shadow\_soundness\_exact.py}.
\end{proposition}
\begin{remark}[Correction to v1]\label{rem:v2fix}
Version~1 of this paper stated the gap backwards. It claimed that a minimal odd-$Q$ support
``can carry the $\Z_2$ shadow obstruction without closing over $\Z_d$'', cited ``an explicit false
certificate at $d=8$ \dots\ recorded in the repository'', and reported that $28.6\%$ of random
$d=8$ incidence systems carrying a $\Z_2$-odd cycle are not genuinely $\Z_d$-contextual. All three
are wrong. The false-positive rate is not $28.6\%$ but exactly $0$, by the proof above; no such
certificate exists to be recorded, which is why the cited artifact was never committed; and the
true defect is the missed-detection direction described above. The error made the criterion look
weaker than it is. What survives is the sentence that followed it, on certificate minimization:
a minimal odd-$Q$ \emph{support} need not be a $\Z_d$-cycle, so it can be a false \emph{witness}
even though the \emph{verdict} is always right. Version~1 conflated witness with verdict.
\end{remark}
\noindent\emph{Why $Q$ is nonetheless a good criterion: a Weyl-specific fact.} On generic $\Z_d$
scenarios the mod-$2$ shadow is a poor proxy --- not because it lies, which it cannot, but because
it is often blind: under the protocol of \texttt{shadow\_soundness\_exact.py} it misses roughly
$39\%$ of genuinely $\Z_8$-contextual random incidence systems. On Weyl families, by contrast, the
$\Z_2$ shadow and the true $\Z_d$ obstruction agree generically ($300/300$ on random families); the
discrepancy appears only on engineered minimal cycle supports. Thus the completeness of the odd-$Q$
criterion is not linear algebra but a consequence of the symplectic structure of the Weyl group, which
forces the carry shadow to be faithful. For certificate minimization one must still search over
$\Z_d$-cycles, since minimality is where the shadow can fail.

\section{The 2-adic tower (structural)}
Beyond the classification lies the question of how the obstruction is generated across dimensions. The
explicit certificates were all found by a uniform construction --- the pure fiber pool over the
minimal attaining family one level below (all local lifts of every base context) --- and the
obstruction propagates up the $2$-adic tower $\mathrm{PM}\to d{=}4\to d{=}8\to d{=}16$ through fibers,
never through sections (Theorem~\ref{thm:depth}). We record what is proved and what remains.
\begin{theorem}[Fiber projection]\label{thm:fiber}
For $\tilde\lambda$ in the mod-$2d$ kernel of a fiber pool $\mathrm{Fib}(F)$, the projection
$\pi(\tilde\lambda)_j=\sum_x\tilde\lambda_{j,x}$ satisfies $\pi^{\mathsf T}A_{\mathrm{base}}\equiv
0\pmod{2d}$ and is everywhere even; hence $\mu=\pi/2$ is a mod-$d$ base kernel element. Consequently
the value group of $\mathrm{Fib}(F)$ at $2d$ is $\subseteq 2\cdot(\text{value group of }F)$ under
$\Z_d\hookrightarrow\Z_{2d}$.
\end{theorem}
\begin{conjecture}[Fiber doubling law, $m_0\ge 4$]\label{conj:doubling}
For a base family $F$ at even level $m_0\ge 4$ outside a generative subclass, the value-group
generators obey $g_{\mathrm{fiber}}=2\,g_{\mathrm{base}}$. The reflect direction ($F$ attains
$\Rightarrow\mathrm{Fib}(F)$ attains) holds on every attaining family tested; the preserve direction
($F$ trivial $\Rightarrow\mathrm{Fib}(F)$ trivial) fails on a positive fraction of trivial bases even
at $m_0=4$ (see below), so the law holds exactly on the non-generative subclass, characterized in
closed form below.
\end{conjecture}
\noindent\emph{Status.} Verified $51/51$ across attaining anchors (PM, cert4, $\mathrm{ROZF}(d)$,
$d\le 10$) and random families, and $10/10$ on ROZF families $d=4..22$. The reflect direction is solid
on every attaining family tested. Preserve, however, leaks: pre-filtering random bases by actual
base-triviality, a positive fraction (empirically ${\sim}12\text{--}15\%$ at $m_0=4,6$) of trivial
bases have an attaining fiber. So the remaining obligation is not a blanket $\delta$-vanishing lemma
but a characterization of the non-generative trivial bases on which every fiber certificate has value
$0$. Via Theorems~\ref{thm:valuebit},~\ref{thm:fiber} this is a finite $\mathrm{GF}(2)$ statement, and
this session localizes the obstruction precisely to the carry of $\tfrac12(T_A+T_B)$ in the value-bit
formula (the base-commutator piece $T_A$ and the high quadratic $(d/2)T_Q$ both vanish for
$4\mid m_0$; integer halving does not distribute over $T_A+T_B$, and that carry is the entire
content). [open]

\paragraph{The generativity mechanism, made exact.}
Using the criterion of Theorem~\ref{thm:cycle}, the fiber cycle value decomposes exactly. Writing each
fiber observable as $\tilde v=u+d\,x$ ($u\in\Z_d$ the base residue, $x\in\{0,1\}$ the lift bit), the
fiber $Q$ is
\[
Q_{\mathrm{fiber}}=Q_{\text{base-residue}}\ \oplus\ \sum_{a<b}\Big\lfloor
\tfrac{\ip{u_a}{u_b}+d\,M_{ab}+d^2\ip{x_a}{x_b}}{2d}\Big\rfloor-\Big\lfloor\tfrac{\ip{u_a}{u_b}}{2d}
\Big\rfloor\pmod 2,
\]
with $M_{ab}=\ip{u_a}{x_b}+\ip{x_a}{u_b}$ (the underbraced term is the lift-induced halving carry).
Generativity is precisely the event that this lift-induced carry is odd while the base is trivial.
Crucially it is a floor-carry, not a quadratic form in the lift bits --- which is why an Arf/quadratic
treatment of the doubling law does not close it: the obstruction is of the same carry type as the
value bit $b$ itself. This carry now has a closed form on all fiber cycles. Writing
$\Sigma M=\sum_{a<b}M_{ab}$ and $K=\#\{a<b:M_{ab}\ \text{odd}\}$,
\[
G=\big\lfloor\tfrac{\Sigma M}{2}\big\rfloor\oplus\big\lfloor\tfrac{K}{2}\big\rfloor\pmod 2,
\]
verified exactly ($7500/7500$ at $d=4$, $1200/1200$ at $d=6$). The first term is the intrinsic pairing
$\Sigma M=\ip{U}{X}-\sum_i\ip{u_i}{x_i}$ of the cycle's total base-residue $U=\sum u_i$ and total
lift-bit $X=\sum x_i$; the second is the carry produced by summing the odd cross-terms (the
sum-of-floors versus floor-of-sum defect), which vanishes exactly when $K$ is even --- in particular
on repeat-free cycles, recovering the intrinsic criterion $\ip{U}{X}-\sum_i\ip{u_i}{x_i}\equiv 0,1
\pmod 4$. The generativity/doubling-law obstruction is thus fully characterized: a cycle is
non-generative iff $G=0$.

The carry decomposition is provable as an exact integer identity. For any integers
$M_1,\dots,M_N$ with $S=\sum_a M_a$ and $K=\#\{a:M_a\ \text{odd}\}$,
\[
\sum_a\big\lfloor\tfrac{M_a}{2}\big\rfloor=\big\lfloor\tfrac{S}{2}\big\rfloor\oplus
\big\lfloor\tfrac{K}{2}\big\rfloor\pmod 2,
\]
since $\lfloor M/2\rfloor=(M-(M\bmod 2))/2$ gives $\sum_a\lfloor M_a/2\rfloor=(S-K)/2$, while
$S\equiv K\pmod 2$ gives $\lfloor S/2\rfloor=(S-(K\bmod 2))/2$, and the difference is
$-\lfloor K/2\rfloor$. Hence the closed form for $G$ is proved once $Q_{\mathrm{fiber}}$ is the sum of
the per-pair halving carries $\lfloor M_{ab}/2\rfloor$ (the one remaining verified geometric step); the
arithmetic that turns the per-pair carries into the two-term formula is exact.

\begin{remark}[Generativity and contextuality activation]\label{rem:activation}
The fiber can be generative: a base with trivial value group can have an attaining fiber, creating
contextuality absent below. This is strongest at $m_0=2$ but, correcting an earlier expectation,
persists at $m_0=4,6$ at a ${\sim}12\text{--}15\%$ rate per random trivial base; generativity is a
generic tower phenomenon, not a qubit-only artifact. Read as a protocol, this is contextuality
activation by dimension embedding: a passive embedding $\Z_d\hookrightarrow\Z_{2d}$ (a free operation
--- no entanglement or magic) converts a noncontextual family into a state-independent contextual one.
The yield is tunable by the number of base contexts $k$, rising monotonically $5.6\%$ ($k{=}3$)
$\to 76.4\%$ ($k{=}12$) ($250$ trials per point, Wilson $95\%$ CI), plausibly $\to 1$ for large $k$.
Because the witness is state-independent, the certified nonclassicality holds for any input state,
including the maximally mixed one --- a preparation-assumption-free feature. Crucially, generativity
is a single-context fiber effect and is structurally distinct from the multi-context certificate
cycles that carry the value-$d/2$ obstruction (see Section~\ref{sec:open}).
\end{remark}

\paragraph{Connection to quasiprobability negativity.} The equivalence between (measurement)
noncontextuality and the existence of a nonnegative quasiprobability representation is a general
theorem~\cite{Spekkens,Wallman}; on the quasiprobability side, no nonnegative diagram-preserving
representation exists in any even dimension~\cite{Wallman}, which is the even/odd dichotomy seen here
from the dual direction. Our contribution in this language is constructive and quantitative: the
odd-$Q$ / $\Z_d$-cycle criterion is a closed-form, decidable test for whether a Weyl measurement
family admits such a representation (the general equivalence is non-constructive), and the value
classification refines the binary present/absent statement to a fixed obstruction strength $d/2$. We
add no new simulation bound: the state-independent obstruction places even-$d$ Weyl measurement
families on the negative-representation side of the measurement axis, but classical simulation cost is
governed by state negativity, a distinct quantity.

\section{Related work}
Existence of parity/AvN obstructions and the cohomological criterion ``proof exists $\iff[\beta]\neq
0$'' are due to Okay--Roberts--Bartlett--Raussendorf~\cite{ORBR}; our certificate framework repackages
the existence side. Even-$d$ cohomological obstructions and general-$d$ attainment are ROZF~\cite{ROZF};
odd-$d$ noncontextuality and the $2^{-1}\bmod d$ mechanism are Gross and Raussendorf et
al.~\cite{Gross,RBDOB}. Linear-algebra enumeration of $\Z_2$ parity proofs (product $\pm I$ only, no $\Z_d$
spectra, no towers) is Lison\v{e}k--Raussendorf--Singh~\cite{LRS}. The formulation of observable-based
contextuality as the unsolvability of a $\mathrm{GF}(2)$ linear system, with the qubit Pauli group
encoded in the symplectic polar space $W_n$ and a numerically computed degree of contextuality
(minimum unsatisfied weight), is the Saniga--Holweck--de~Boutray--Giorgetti
program~\cite{deBoutray,Muller}; it is qubit-only ($d=2$) and computes the decision by SAT/heuristic
search. Our even-$d$ commutator-carry $b$ and incidence matrix $A$ are the $\Z_d$ lift of their $W_n$
encoding; the distinction is that our obstruction admits a closed-form per-cycle evaluation
(Theorem~\ref{thm:cycle}) rather than a min-weight search, and extends to even qudits with a value
spectrum. Conversely, their quantitative degree of contextuality $d_H(E,\mathrm{Im}\,A)$ lifts verbatim
to even $d$ through our carry vector $b$: it counts how many context constraints must be flipped to
restore a valuation, and is a distinct, orthogonal axis to the value classification (a family of
degree $1$ can carry the maximal value $d/2$). A qudit contextuality-degree over $\Z_d$ (rather than
the $\Z_2$ coset weight) is a natural open refinement. A complete characterization of which qubit
stabilizer states admit an AvN proof (the AvN-triple theorem: every such argument reduces to a
GHZ-equivalent triple) is Abramsky--Barbosa--Car\`u--Perdrix~\cite{ABCP}, with completeness /
false-negative analysis by Car\`u and Abramsky et al.~\cite{ABKLM}; a state-independent AvN class via
partial closure of commuting observables is Kim--Abramsky~\cite{KimAbramsky}. \textbf{The even/odd dichotomy and the value classification $\{0,d/2\}$ are due to Abramsky--Cercelescu--Constantin~\cite{ACC24}} (parity obstruction $2k\equiv0$, their Thm.~16; commutation-group normal-form classification, their Thm.~21), who also give a $\Z_2\to\Z_{2k}$ transfer construction. Our Theorem~\ref{thm:spectrum} recovers their result in Weyl-certificate language; the contribution here is the explicit carry/value-bit calculus, the $2$-adic tower (Depth Rigidity and the doubling law), and the hardware, with the geometric/cohomological account of the obstruction in the companion paper. All of this AvN
literature works over $\Z_2$ (qubit stabilizers, product $\pm I$) and characterizes AvN existence. Our
distinct contributions are of a different kind: the even-$d$ value classification $H(d,m)=\{0,d/2\}$ (a
spectrum over $\Z_d$, not existence over $\Z_2$); the closed commutator-carry form
(Theorem~\ref{thm:carry}); Depth Rigidity with a certificate-derived no-lift witness
(Theorem~\ref{thm:depth}); the explicit high-$d$ certificates; and the fiber-tower program.

\paragraph{Relation to Abramsky--Cercelescu--Constantin: dictionary and logical status.}
Under the identifications
\begin{center}\small
\begin{tabular}{ll}
\hline
\emph{this paper (Weyl realization)} & \emph{Abramsky--Cercelescu--Constantin~\cite{ACC24}}\\
\hline
Weyl label $v$ & commutation-group generator / vector\\
symplectic form $\ip{u}{v}$ & commutator matrix entry $\mu(x,y)=v_{x,y}$\\
closed context $C$ & commuting subword / relation\\
certificate $\lambda$ & contextual word\\
value $S$ & central phase $k$\\
carry $b(C)$ & commutation factor $\sum\mathcal I(w)$ (inversion sum)\\
$\Z_d$ value assignment & left splitting / noncontextual homomorphism\\
\hline
\end{tabular}
\end{center}
Theorem~\ref{thm:spectrum} is the Weyl--Heisenberg \emph{specialization} of their Theorem~16: their
inversion sum $\sum\mathcal I(w)$ is our carry accumulation, their polarization
$2\sum\mathcal I(w)=\sum_{x<y}n_xn_y\,v_{yx}$ is our polarization of $\beta$, and their $2k\equiv0\pmod d$
is our $2S\equiv0\pmod d$. Our spectrum theorem is therefore neither independent of nor a strengthening
of theirs; it re-derives their classification in the Weyl realization, and the carry $b(C)$ is their
commutation factor in symplectic-form coordinates --- a change of realization and an explicit closed
form, not a new invariant. What is \emph{not} present in~\cite{ACC24} and is offered here as the
candidate new content: the explicit value-bit closed form for $2$-adically lifted certificates
(Theorem~\ref{thm:valuebit}); the \emph{cross-dimensional} tower statements --- Depth Rigidity
(Theorem~\ref{thm:depth}) and the doubling law --- which concern relations \emph{between} $d$ and $2d$
rather than the parity constraint at a single $d$ (and must be read as such, not as restatements of it);
the hardware replication; and the geometric/cohomological home (companion paper). A full
theorem-by-theorem equivalence with~\cite{ACC24}'s normal-form invariants is the natural next step.
The extended arXiv version of~\cite{ACC24} additionally provides a Darboux-normal-form reduction and
a closed-form relative-parity criterion for the existence of state-independent contextuality in
general commutation groups; that existence classification is complementary to the value spectrum and
state-sector facet geometry developed here. Information-theoretic ``cost of contextuality'' measures
\cite{BookkeepingCost,CoordinationCost} quantify simulation overhead rather than the cohomological
obstruction value studied here.

\section{Open problems}\label{sec:open}
(1) The $\delta$-vanishing lemma, closing Conjecture~\ref{conj:doubling} for $m_0\ge 4$. The
closed-form criterion (Theorem~\ref{thm:cycle}) settles the static contextuality question exactly ---
contextual iff some cycle has odd $Q$ --- and clarifies the tower question: propagation $d\to 2d$ is
now the statement that the odd-$Q$ property transfers from base cycles to fiber cycles. The certificate
quadratic form $q$ is genuine (third differences vanish, verified) and on a multi-cycle space the
fiber-attainment map is quadratic with nonzero polarization, so an Arf invariant governs the fiber
transfer. We caution that an earlier ``multi-cycle $\Rightarrow$ contextual'' expectation is false: by
Theorem~\ref{thm:cycle}, generic cycles have $Q=0$ and are noncontextual (verified $300/300$ on random
$m{=}2$ families); contextuality requires the special odd-$Q$ condition. (2) The activation subclass of
Remark~\ref{rem:activation} as a protocol (the generativity carry $G$ is now closed; what remains is
its operational reach). (3) Ramification necessity: do attaining families always live on the mod-$d$
colliding locus? (both known minimal attaining families do). (4) Minimize the $d=16$ certificate.
(5) AvNR/Car\`u-style completeness of the admissible-cycle certificates.

\paragraph{The geometric step of the doubling law.}
We record the identity $Q_{\mathrm{fiber}}\equiv\sum_{a<b}\lfloor M_{ab}/2\rfloor\pmod 2$ together with
its proof, as it completes the value-transfer arithmetic. Writing each level-$2d$ observable as
$v=u+d\,x$ ($u$ the base representative mod $d$, $x\in\{0,1\}^{2m}$ the lift bits) and
$q_u=\lfloor\ip{u_a}{u_b}/d\rfloor$, $M_{ab}=\ip{u_a}{x_b}+\ip{x_a}{u_b}$, $s_{xx}=\ip{x_a}{x_b}$, the
exact per-pair carry is
\[
\Big\lfloor\tfrac{\ip{v_a}{v_b}}{2d}\Big\rfloor\equiv\Big\lfloor\tfrac{q_u+M_{ab}+d\,s_{xx}}{2}
\Big\rfloor\pmod 2
\]
(a pure floor identity, since $0\le\ip{u_a}{u_b}\bmod d<d$ prevents the remainder from carrying;
verified exactly at $d=4,6,8$). For even $d$ this splits as
$\lfloor q_u/2\rfloor+\lfloor M_{ab}/2\rfloor+\tfrac d2 s_{xx}+(q_u\bmod 2)(M_{ab}\bmod 2)$, so
\[
Q_{\mathrm{fiber}}-\sum_{a<b}\big\lfloor\tfrac{M_{ab}}{2}\big\rfloor\equiv\sum_{a<b}
\big\lfloor\tfrac{q_u}{2}\big\rfloor+\tfrac d2 s_{xx}+(q_u\bmod 2)(M_{ab}\bmod 2)\pmod 2.
\]
The first term $\sum\lfloor q_u/2\rfloor$ depends only on the base $u$'s and is identically zero as a
base-context vector (checked for the minimal attaining families at $d=4,6,8$), so it vanishes on every
fiber. The term $\tfrac d2\sum_{a<b}s_{xx}$ closes in general even $d$: it is nonzero only for
$d\equiv 2\pmod 4$, and there the polarization identity for the quadratic refinement
$Q(x)=\sum_i x_{2i}x_{2i+1}$ gives $\sum_{a<b}\ip{x_a}{x_b}=Q(r_C)+\sum_{v\in C}Q(x_v)$ (using the
fiber closure $\sum_{v\in C}x_v=r_C$, $r_C=\tfrac1d\sum_{v\in C}u_v\bmod 2$), so the term is the
explicit per-observable coboundary $\phi(v)=\tfrac d2 Q(x_v)$ plus the base-context constant
$\tfrac d2 Q(r_C)$, itself a base coboundary (verified $d=4,6,10,14$). The last term
$\sum_{a<b}(q_u\bmod 2)(M_{ab}\bmod 2)$ reorganizes as $\sum_v\ip{w_v}{x_v}$ with
$w_c=\sum_{a\neq c}(q_{u_a}\bmod 2)\,u_a$: a per-observable coboundary of linear form in the lift bits,
whose $\F_2$-system is solvable with no base-context correction ($t\equiv 0$), verified at
$d=4,6,8,10,12,14$; hence this term too lies in the column space of $A$ and pairs to zero with every
cycle. This solvability holds for all even $d$, which completes the geometric step. Split by $d\bmod 4$:
for $d\equiv 0\pmod 4$ the base symbol $u=\tfrac d2 p$ has $\tfrac d2$ even, hence $u\equiv 0\pmod 2$
componentwise; since $M_{ab}$ is bilinear, $M_{ab}\equiv 0\pmod 2$ on every fiber context, so the term
$\equiv 0$ identically. For $d\equiv 2\pmod 4$, $\tfrac d2$ is odd, so $u\equiv p\pmod 2$ and
$M_{ab}\bmod 2$ depends only on the fixed base symbols $p$ and the lift bits $x$, not on $d$; labeling
observables and contexts by their $d$-independent (symbol, lift) signature, the coboundary matrix $A$
and the target vector $T_{\mathrm{mix}}$ are identical for every $d\equiv 2\pmod 4$ --- one fixed
$768\times 144$ $\F_2$ system (byte-identical for $d=6,10,14,18,22,26,30$, weight
$|T_{\mathrm{mix}}|=128$). Solvability is therefore a single fact about this fixed system, established
at $d=6$, and transfers verbatim to every $d\equiv 2\pmod 4$.
The complementary class $d\equiv 0\pmod 4$ is settled separately, and more easily: writing
$u=hp \bmod d$ with $h=d/2$, bilinearity gives $q_u=\mathrm{symp}(u_a,u_b)/d=h s$ where
$\mathrm{symp}(p_a,p_b)=2s$, so $h$ even forces every $q_u$ to be even and $T_{\mathrm{mix}}$
to vanish \emph{identically}, pointwise, with no cocycle summation required. The whole system is
thus a function of $d\bmod 4$ alone, constant within each class and vanishing outright on one of
them (verified $d=4,\dots,20$; see \texttt{mod4\_dichotomy.py}).
Hence $T_{\mathrm{mix}}$ is a coboundary for all even $d$; together with the vanishing of $\sum\lfloor q_u/2\rfloor$ and the explicit potential
$\phi(v)=\tfrac d2 Q(x_v)$ for the $s_{xx}$ term, the geometric step of the doubling law is established
for all even $d$.

\section{Hardware validation at the base of the tower ($d=2$)}
The $d=2$ endpoint of the spectrum---value $d/2=1$, i.e.\ the Peres--Mermin obstruction---was executed
on a superconducting device (IBM \texttt{ibm\_marrakesh}) to validate the protocol of the preceding
remark. Each of the six contexts was realized as a three-observable sequential non-destructive
measurement (one shared ancilla, Hadamard-test readout of each commuting Pauli into its own bit, with
reset between), so that the per-context correlator is not algebraically pinned and genuinely responds
to device error. The Cabello witness
$S=\langle R_1\rangle+\langle R_2\rangle+\langle R_3\rangle+\langle C_1\rangle+\langle C_2\rangle
-\langle C_3\rangle$ has noncontextual bound $4$ and quantum value $6$. On each device we pinned all
six contexts to a single low-error qubit triple (selected by readout- and two-qubit-gate error),
enabled dynamical decoupling and measurement twirling, and took $8192$ shots per context. Replicating
the identical protocol across three IBM Heron devices gave, per device,
\[
S_{\mathrm{marrakesh}}=4.94\pm 0.02,\quad S_{\mathrm{fez}}=4.56\pm 0.02,\quad
S_{\mathrm{kingston}}=4.74\pm 0.02,
\]
each individually violating the bound by $32$--$61\sigma$, with an inverse-variance-weighted
combination
\[
S=4.761\pm 0.010\qquad(80\sigma\text{, shot-noise only}),
\]
and per-context magnitudes $0.74$--$0.85$ (e.g.\ $\langle C_3\rangle$ ranging $-0.74$ to $-0.82$),
consistent with reported superconducting realizations. This combined figure is shot-noise-only: the
three devices are \emph{not} mutually consistent (a homogeneity test gives $\chi^2\approx266$ on $2$
degrees of freedom), so once the chip-to-chip scatter is honored a systematics-aware (Birge-scaled)
combination gives $\approx7\sigma$. The robust statement is therefore the \emph{per-device} violation
($32$--$61\sigma$, shot noise); the $80\sigma$ combined number should not be read as a single-$S$
significance, and no nondisturbance/compatibility diagnostic has been performed (repository
\texttt{hardware/COMPATIBILITY.md}). The $0.38$ spread in $S$ across devices is
genuine chip-to-chip variation and leaves every device far above the classical bound. An initial run
in which contexts were transpiled independently placed one context on poorer qubits and yielded a
marginal $S=4.06\pm 0.03$ ($2.0\sigma$); the fixed-layout, error-suppressed, multi-device figure above
is the reliable one. We report this only as a validation of the measurement methodology at the base of
the tower; it certifies the presence of the $d=2$ obstruction, not the value grading, whose hardware
realization at $d=4$ we leave to the separate value-graded study. We note that, by
Corollary~\ref{cor:bit}, the obstruction phase equals $e^{2\pi i(d/2)/d}=-1$ at every even $d$, so
higher-$d$ hardware witnesses exhibit the same $-1$ signature: the $d$-dependence of the spectrum is an
algebraic ($\Z_d$) datum, certified in software, rather than a distinct analog phase on-device.

\appendix
\section{Computational certificates}
All load-bearing certificates ship in the companion repository, with each claim mapped to a script and
its pinned one-line output in \texttt{INDEX.md}:
\texttt{spectrum\_test2.py} (potential identity, odd-$d$ Gross convention);
\texttt{Wformula.py} (Theorem~\ref{thm:valuebit});
\texttt{verify\_cert8.py}, \texttt{verify\_cert16.py} (raw-matrix $d=8,16$ certificates);
\texttt{criterion.py} (the closed-form odd-$Q$ criterion);
\texttt{tmix\_dindep.py}, \texttt{close\_T2\_proof.py} (the fixed $\F_2$ system and the geometric step
of the doubling law); \texttt{kernelpk.py}, \texttt{howell.py} (Smith/Howell-form solvers,
brute-force unit-tested). The additional stress tests reported inline (large-sample potential-identity
and carry-form scans, the ROZF attainment sweep $d\le 20$, the fiber doubling-law statistics, and the
$d=8$ false-certificate example) are archived in the same repository ledger.

\begin{thebibliography}{99}
\bibitem{ORBR} C.~Okay, S.~Roberts, S.~D.~Bartlett, R.~Raussendorf,
\emph{Quantum Inf.\ Comput.} \textbf{17}, 1135 (2017), arXiv:1701.01888.
\bibitem{ROZF} R.~Raussendorf, C.~Okay, M.~Zurel, P.~Feldmann,
\emph{Quantum} \textbf{7}, 979 (2023), arXiv:2110.11631.
\bibitem{Gross} D.~Gross, \emph{Hudson's theorem for finite-dimensional quantum systems},
J.~Math.\ Phys.\ \textbf{47}, 122107 (2006), arXiv:quant-ph/0602001.
\bibitem{RBDOB} R.~Raussendorf, D.~E.~Browne, N.~Delfosse, C.~Okay, J.~Bermejo-Vega,
\emph{Contextuality and Wigner function negativity in qubit quantum computation},
\emph{Phys.\ Rev.\ A} \textbf{95}, 052334 (2017), arXiv:1511.08506.
\bibitem{LRS} P.~Lison\v{e}k, R.~Raussendorf, V.~Singh, arXiv:1401.3035.
\bibitem{Spekkens} R.~W.~Spekkens, \emph{Negativity and contextuality are equivalent notions of
nonclassicality}, \emph{Phys.\ Rev.\ Lett.} \textbf{101}, 020401 (2008).
\bibitem{Kirchmair} G.~Kirchmair et al., \emph{State-independent experimental test of quantum
contextuality}, \emph{Nature} \textbf{460}, 494 (2009).
\bibitem{Qu} D.~Qu, K.~Wang, L.~Xiao, X.~Zhan, P.~Xue, \emph{State-independent test of quantum
contextuality with either single photons or coherent light}, \emph{npj Quantum Inf.} \textbf{7}, 154
(2021).
\bibitem{EJP} \emph{A demonstration of contextuality using quantum computers}, \emph{Eur.\ J.\ Phys.}
\textbf{43}, 045403 (2022), DOI 10.1088/1361-6404/ac79e0 (Peres--Mermin square on IBM superconducting
qubits).
\bibitem{Wallman} J.~J.~Wallman, S.~D.~Bartlett, \emph{Non-negative subtheories and quasiprobability
representations of qubits}, \emph{Phys.\ Rev.\ A} \textbf{85}, 062121 (2012), arXiv:1203.2652.
\bibitem{deBoutray} H.~de~Boutray, F.~Holweck, A.~Giorgetti, P.-A.~Masson, M.~Saniga,
\emph{Contextuality degree of quadrics in multi-qubit symplectic polar spaces}, \emph{J.\ Phys.\ A}
\textbf{55}, 475301 (2022), arXiv:2105.13798.
\bibitem{Muller} A.~Muller, M.~Saniga, A.~Giorgetti, F.~Holweck, C.~Kelleher, \emph{A new heuristic
approach for contextuality degree estimates}, arXiv:2407.02928.
\bibitem{KimAbramsky} B.~Kim, S.~Abramsky, \emph{State-independent all-versus-nothing arguments},
arXiv:2311.11218 (2023).
\bibitem{ACC24} S.~Abramsky, \c{S}.-I.~Cercelescu, C.~Constantin, \emph{Commutation Groups and State-Independent Contextuality}, FSCD 2024, LIPIcs \textbf{299}, 28:1--28:20; extended version arXiv:2603.12197 (2026).
\bibitem{BookkeepingCost} arXiv:2601.20167 (2026).
\bibitem{CoordinationCost} arXiv:2606.23577 (2026).
\bibitem{ABCP} S.~Abramsky, R.~Soares~Barbosa, G.~Car\`u, S.~Perdrix, \emph{A complete characterisation
of All-versus-Nothing arguments for stabiliser states}, \emph{Phil.\ Trans.\ R.\ Soc.\ A} \textbf{375},
20160385 (2017), arXiv:1705.08459.
\bibitem{ABKLM} S.~Abramsky, R.~Soares~Barbosa, K.~Kishida, R.~Lal, S.~Mansfield, \emph{Contextuality,
cohomology and paradox}, CSL 2015; G.~Car\`u, false-negative analysis (2018).
\end{thebibliography}

\end{document}
